\section{Результаты}

Была разработана программа, реализующая межпроцессное взаимодействие по схеме «родитель-потомок» с использованием двух неименованных каналов (pipe) в операционной системе Linux.

{\bfseries Ключевые особенности решения:}

1. Программа использует интерфейс `os.h`, который инкапсулирует системные вызовы, что теоретически позволяет легко портировать код на другие платформы.

2. Реализовано два канала:
   - `pipe1`: передача строк от родителя к потомку
   - `pipe2`: передача сообщений об ошибках от потомка к родителю

3. Все файловые дескрипторы и каналы закрываются своевременно, предотвращая утечки ресурсов и зависания программы.

4. Программа проверяет возвращаемые значения всех системных вызовов и выводит понятные сообщения об ошибках.

5. Родительский и дочерний процессы представлены разными исполняемыми файлами, что соответствует требованиям задания.